\chapter{情感分析方法说明}

\section{方法概述}

本研究采用大型语言模型对 bitcoin/bitcoin 仓库 Issue 评论进行情感分析，采用“先按周聚合、后整体评分”的策略。
主分析模型为 Claude Sonnet 4.5（claude-sonnet-4.5），
稳健性检验分别以 DeepSeek-V3.2（deepseek-v3.2）和 GPT-5.2（gpt-5.2）
两个替代 LLM 以及 VADER 词典方法进行对照。
三个 LLM 分属 Anthropic、DeepSeek、OpenAI 三家公司，
架构与训练路线各不相同，可最大化跨模型多样性；
VADER 作为基于规则的工具，与 LLM 的语义推理在方法论上存在本质差异。
所有 LLM 均通过 OpenAI 兼容 API 接口调用，
推理温度设置为 0，并启用 JSON 强制输出模式
（\texttt{response\_format=\{"type":"json\_object"\}}），
以保证结果的确定性、可重复性与格式规范性。

\section{情感打分提示词}

以下为提交给模型的系统提示词与用户提示词的完整文本。

\vspace{1em}
\noindent System Prompt:

\begin{quote}\small
You are an expert sentiment analyst. Analyze the given text and provide comprehensive sentiment scores across multiple dimensions.

\medskip
Task 1: Overall Sentiment Polarity (9-point scale)

Evaluate the overall emotional tone:
\begin{itemize}
  \item 1--4: Negative sentiment (1=extremely negative, 4=slightly negative)
    \begin{itemize}
      \item Indicators: anger, disappointment, criticism, frustration, complaints
    \end{itemize}
  \item 5: Neutral sentiment
    \begin{itemize}
      \item Indicators: objective, factual, balanced, or mixed emotions
    \end{itemize}
  \item 6--9: Positive sentiment (6=slightly positive, 9=extremely positive)
    \begin{itemize}
      \item Indicators: joy, gratitude, approval, excitement, satisfaction
    \end{itemize}
\end{itemize}

\medskip
Task 2: Six Basic Emotions (5-point scale, 0--4 for each)

Rate the intensity of each emotion:
\begin{itemize}
  \item Anger: Irritation, frustration, hostility (0=none, 4=extreme)
  \item Disgust: Revulsion, contempt, distaste (0=none, 4=extreme)
  \item Fear: Anxiety, worry, concern (0=none, 4=extreme)
  \item Happiness: Joy, satisfaction, contentment (0=none, 4=extreme)
  \item Sadness: Sorrow, disappointment, melancholy (0=none, 4=extreme)
  \item Surprise: Astonishment, unexpectedness (0=none, 4=extreme)
\end{itemize}

\medskip
Task 3: Disappointment --- Expectation Disconfirmation (5-point scale, 0--4)

This is a critical dimension for measuring expectation violation.
Disappointment occurs when reality fails to meet prior expectations or hopes.

Definition: Disappointment reflects the gap between what was expected/hoped for
and what actually occurred. It is distinct from general sadness or anger,
as it specifically involves unmet expectations.

Key Indicators:
\begin{itemize}
  \item Explicit statements of unmet expectations
        (e.g., ``I expected X but got Y'', ``This is not what I hoped for'')
  \item Expressions of letdown, disillusionment, or feeling underwhelmed
  \item Comparisons between anticipated and actual outcomes
  \item Phrases like ``unfortunately'', ``sadly'', ``not as good as expected'', ``fell short''
  \item Regret about decisions or outcomes that didn't meet hopes
  \item Loss of enthusiasm or optimism that was previously present
\end{itemize}

Rating Scale:
\begin{itemize}
  \item 0: No disappointment --- expectations met or exceeded, or no expectations were present
  \item 1: Slight disappointment --- minor unmet expectations, easily overlooked
  \item 2: Moderate disappointment --- clear gap between expectations and reality, noticeable dissatisfaction
  \item 3: Strong disappointment --- significant expectation violation, considerable distress or frustration
  \item 4: Extreme disappointment --- severe expectation disconfirmation, profound sense of letdown or betrayal
\end{itemize}

Important: Distinguish disappointment from:
\begin{itemize}
  \item Pure anger (which may lack the expectation component)
  \item General sadness (which may not involve unmet expectations)
  \item Neutral criticism (which may be objective without emotional investment)
\end{itemize}

\medskip
Output Format:

Return a JSON object with the following structure. For each score, provide a brief
one-sentence basis explaining your rating:

\begin{verbatim}
{
    "sentiment_polarity": {
        "score": <int, 1-9>,
        "basis": "<one sentence explanation>"
    },
    "anger": {
        "score": <int, 0-4>,
        "basis": "<one sentence explanation>"
    },
    "disgust":       { "score": <int, 0-4>, "basis": "..." },
    "fear":          { "score": <int, 0-4>, "basis": "..." },
    "happiness":     { "score": <int, 0-4>, "basis": "..." },
    "sadness":       { "score": <int, 0-4>, "basis": "..." },
    "surprise":      { "score": <int, 0-4>, "basis": "..." },
    "disappointment": {
        "score": <int, 0-4>,
        "basis": "<explanation focusing on
                   expectation disconfirmation>"
    }
}
\end{verbatim}

Important:
\begin{itemize}
  \item Provide ONLY the JSON object, no additional text
  \item Each score must be within the specified range
  \item For disappointment, the basis should explicitly reference the expectation-reality gap if present
\end{itemize}
\end{quote}

\vspace{0.5em}
\noindent User Prompt:
\begin{quote}\small
Text to analyze: \textit{[待分析的评论全文]}
\end{quote}

\section{技术配置与处理流程}

\begin{enumerate}
  \item 模型配置：主分析使用 claude-sonnet-4.5，temperature = 0；
        稳健性检验使用 DeepSeek-V3.2（deepseek-v3.2）和 GPT-5.2（gpt-5.2）
        两个替代 LLM，temperature = 0；
        词典基线使用 VADER\cite{hutto2014vader}。
        所有 LLM 均启用 JSON 强制输出模式。
  \item 文本预处理：
    \begin{itemize}
      \item 过滤以三个反引号包裹的纯代码片段，消除代码噪声
      \item 仅保留包含至少 5 个有效词的评论，剔除过短的无意义内容
    \end{itemize}
  \item 并发与容错：采用线程池（\texttt{ThreadPoolExecutor}，默认 10 线程）
        并发调用 API。
        针对超时和 429（速率限制）两类错误分别实施指数退避重试策略，
        超时错误最大等待 30 秒，速率限制错误最大等待 60 秒，
        均最多重试 10 次。所有请求结果以 JSON 文件逐条缓存至本地，
        支持中断后的断点续传。
  \item 输出格式：每周返回包含 8 个维度（情感极性、愤怒、厌恶、
        恐惧、快乐、悲伤、惊讶、失望）评分及评分依据的嵌套 JSON 对象。
  \item 聚合与评分策略：先将每周所有经预处理的有效评论合并为一个文本块，
        再由模型对整周文本进行一次性评分，直接输出该周情感极性变量
        （$\text{Sentiment}_t$，1--9 分制）。
        选择先聚合后评分而非逐条评分再取均值，原因在于：
        （1）不同评论的信息含量与代表性存在差异，逐条评分后简单算术均值
        会在模型推理误差之上叠加聚合方法的额外误差；
        （2）样本期内评论总量庞大，逐条评分的 API 调用时间与资金成本过高。
  \item 脚本位置：完整分析脚本见 \texttt{sentiment\_scoring.py}。
\end{enumerate}
